\SourceLeaf{364}%
einfach durch \(240:1823\) ausdrücken. Auf diesen Vorteil hätte er verzichten müssen, wenn er ursprünglich vom Bruch \(\frac{362}{209}\) ausgegangen wäre.

Am Ende der Abhandlung bemerkt Herr M. noch, dass auch \textsc{Columella} in \textit{de re rustica} V, 2 einen der Näherungswerte von \(\BareRoot{3}\), nämlich \(\frac{26}{15}\), verwendet hat. Für den Flächeninhalt des gleichseitigen Dreiecks setzt er die Summe eines Drittels und eines Zehntels des auf seiner Seite errichteten Quadrats an.

